\chapter{\'Etat de l'Art et Choix Technologiques}
\minitoc


%--------------------------------------------------------------
\section{Introduction}
%--------------------------------------------------------------

\todo{Annoncer la structure : 9+ couches technologiques examinées,
chacune avec comparaison et décision justifiée.}

%--------------------------------------------------------------
\section{Agents conversationnels vocaux}
%--------------------------------------------------------------

\subsection{Taxonomie}

\todo{IVR classique $\rightarrow$ voicebot rule-based $\rightarrow$ voicebot NLU
$\rightarrow$ voicebot LLM-augmenté.
Référence : Orange Innovation TALN 2022 (arXiv 2404.02009).}

\begin{figure}[H]
  \centering
  \input{tikz/fig2_taxonomie_voicebot}
  \caption{Taxonomie des agents conversationnels vocaux}
  \label{fig:taxonomie_voicebot}
\end{figure}

\subsection{Architecture canonique d'un voicebot}

\todo{%
  Présenter l'architecture générique STT $\rightarrow$ NLU $\rightarrow$
  Dialogue Manager $\rightarrow$ TTS avant de détailler les choix technologiques.
  Cette figure sert de référence pour toutes les comparaisons du chapitre.
}

\begin{figure}[H]
  \centering
  \input{tikz/fig2_archi_canonique}
  \caption{Architecture canonique d'un voicebot NLU}
  \label{fig:archi_canonique}
\end{figure}

\todo{Voicebot NLU + LLM-augmenté, 100\% on-premise.}

%--------------------------------------------------------------
\section{Couche téléphonique et traitement audio}
%--------------------------------------------------------------

\subsection{Protocoles SIP/RTP}

\todo{Définition, rôle dans un système d'appels sortants.}

\subsection{Comparaison des serveurs SIP}

\begin{table}[H]
  \centering
  \caption{Comparaison des serveurs SIP}
  \label{tab:comp_sip}
  \begin{tabularx}{\textwidth}{lXXX}
    \toprule
    \textbf{Critère} & \textbf{FreeSWITCH} & \textbf{Asterisk} & \textbf{Twilio} \\
    \midrule
    Open-source         & \checkmark & \checkmark & $\times$ (SaaS) \\
    Auto-hébergement    & \checkmark & \checkmark & $\times$ \\
    mod\_audio\_stream  & \checkmark (natif) & $\times$ & N/A \\
    Performance         & Haute & Moyenne & N/A \\
    Conformité bancaire & \checkmark & \checkmark & $\times$ \\
    \midrule
    \rowcolor{lightGray}
    \textbf{Retenu}     & \textbf{\checkmark} & & \\
    \bottomrule
  \end{tabularx}
\end{table}

\todo{%
  Paragraphe d'analyse (3-5 phrases) : Twilio elimine d'office par la contrainte
  de residence des donnees. Entre FreeSWITCH et Asterisk, argumenter sur le
  support natif de mod\_audio\_stream qui evite un dev custom de streaming
  audio temps-reel -- c'est le critere decisif, pas juste "performance".
}

\subsection{Détection d'activité vocale (VAD)}

\todo{%
  Redaction en prose (PAS de tableau ici -- un seul critere de comparaison
  ne justifie pas un tableau) : expliquer le role du VAD, puis comparer
  Silero VAD V5 a WebRTC VAD en 1 paragraphe (precision, faux positifs,
  integration Python). Conclure sur le choix de Silero VAD V5.
}

%--------------------------------------------------------------
\section{Reconnaissance automatique de la parole (STT)}
%--------------------------------------------------------------

\subsection{Évolution des architectures}

\todo{GMM-HMM $\rightarrow$ CNN-RNN $\rightarrow$ Transformer (Whisper).
Défi linguistique tunisien (arXiv 2309.11327, arXiv 2504.02604).}

\subsection{Comparaison des solutions STT}

\begin{table}[H]
  \centering
  \caption{Comparaison des solutions \ac{STT}}
  \label{tab:comp_stt}
  \begin{tabularx}{\textwidth}{lXllX}
    \toprule
    \textbf{Solution} & \textbf{Type} & \textbf{WER FR} & \textbf{Tunisien} & \textbf{Conformité} \\
    \midrule
    Azure Speech       & Cloud & $\sim$5\% & partiel & $\times$ \\
    Google STT         & Cloud & $\sim$6\% & non & $\times$ \\
    OpenAI Whisper     & Local & $\sim$8\% & faible & \checkmark \\
    faster-Whisper     & Local & $\sim$8\% & faible & \checkmark \\
    \midrule
    \rowcolor{lightGray}
    \textbf{Retenu : faster-Whisper} & \textbf{Local} & & & \textbf{\checkmark} \\
    \bottomrule
  \end{tabularx}
\end{table}

\todo{%
  Paragraphe d'analyse (3-5 phrases) : pourquoi faster-Whisper malgré un WER
  legerement superieur a Azure/Google ? Argumenter sur la contrainte de
  residence des donnees (BNF04), le cout d'inference local vs API payante,
  et la possibilite de fine-tuning sur corpus tunisien (perspective Ch. futur).
  Ne JAMAIS laisser un tableau sans interpretation ecrite qui en tire la conclusion.
}

%--------------------------------------------------------------
\section{Compréhension du langage naturel (NLU) et gestion du dialogue}
%--------------------------------------------------------------

\subsection{Approches NLU}

\todo{Rule-based $\rightarrow$ ML classique (SVM) $\rightarrow$ Transformers $\rightarrow$
frameworks dédiés. Benchmark PFIA 2022.}

\subsection{Comparaison Rasa vs alternatives}

\begin{table}[H]
  \centering
  \caption{Comparaison des frameworks \ac{NLU}}
  \label{tab:comp_nlu}
  \begin{tabularx}{\textwidth}{lXXXX}
    \toprule
    \textbf{Critère} & \textbf{Rasa 3.6} & \textbf{Dialogflow CX} & \textbf{LUIS} & \textbf{Wit.ai} \\
    \midrule
    Open-source            & \checkmark & $\times$ & $\times$ & \checkmark \\
    On-premise obligatoire & \checkmark & $\times$ & $\times$ & $\times$ \\
    Flexibilité actions    & Très haute & Moyenne & Moyenne & Faible \\
    Support FR             & \checkmark & \checkmark & \checkmark & \checkmark \\
    Intégration Python     & \checkmark & Partielle & $\times$ & $\times$ \\
    \midrule
    \rowcolor{lightGray}
    \textbf{Retenu}        & \textbf{\checkmark} & & & \\
    \bottomrule
  \end{tabularx}
\end{table}

\todo{%
  Paragraphe d'analyse : citer le benchmark PFIA 2022 (Schaeffer \& Bouvard) --
  LUIS devance legerement Rasa en pure precision, mais l'argument decisif
  reste l'obligation on-premise. Conclure que Rasa est le seul framework
  NLU enterprise-grade 100\% deployable on-premise (deja mentionne en
  encadre plus loin -- eviter la repetition mot pour mot).
}

\subsection{Génération augmentée par récupération (RAG)}

\todo{%
  Redaction en prose : definir le RAG et justifier pourquoi les politiques
  bancaires ne peuvent pas etre hardcodees dans Rasa (volume, evolutivite,
  precision des citations). Comparer en 1 paragraphe ChromaDB a Pinecone
  (cloud, exclu) et Weaviate (plus lourd a deployer). Comparer egalement
  les embeddings bge-m3 (multilingue, local via Ollama) a text-embedding-ada
  (OpenAI, cloud, exclu). Un petit tableau 2 colonnes est acceptable ICI
  seulement si le texte avant/apres reste substantiel -- sinon rester en prose.
}

%--------------------------------------------------------------
\section{LLM et génération de réponses}
%--------------------------------------------------------------

\todo{Rôle précis (complément Rasa, pas remplacement).
Cloud LLMs exclus. Ollama comparé à Cerebras.
Retenu : Cerebras gpt-oss-120b — argument text-only, aucun PII.}

%--------------------------------------------------------------
\section{Synthèse vocale (TTS)}
%--------------------------------------------------------------

\begin{table}[H]
  \centering
  \caption{Comparaison des solutions \ac{TTS}}
  \label{tab:comp_tts}
  \begin{tabularx}{\textwidth}{lXlllX}
    \toprule
    \textbf{Solution} & \textbf{Type} & \textbf{Lat. CPU} & \textbf{Taille} & \textbf{Noms TN} & \textbf{Conf.} \\
    \midrule
    ElevenLabs              & Cloud & $\sim$800ms & N/A     & bonne & $\times$ \\
    Coqui XTTS v2           & Local & $\sim$3-8s  & $\sim$2GB & bonne & \checkmark \\
    Piper TTS               & Local & $\sim$50ms  & $\sim$200MB & mauvaise & \checkmark \\
    \midrule
    \rowcolor{lightGray}
    \textbf{Piper + dict. phonémique} & \textbf{Local} & \textbf{$\sim$50ms} & \textbf{$\sim$200MB} & \textbf{\checkmark} & \textbf{\checkmark} \\
    \bottomrule
  \end{tabularx}
\end{table}

\todo{%
  Paragraphe d'analyse : Piper seul a une mauvaise prononciation des noms
  tunisiens -- expliquer ICI le probleme IPA via brackets espeak (non fiable),
  puis la solution retenue : dictionnaire orthographique francais
  (tunisian\_phonemes.json, 226 entrees). C'est ce qui transforme Piper
  d'une solution "mauvaise sur les noms TN" a la solution retenue --
  expliciter cette transformation, pas juste l'annoncer.
}

%--------------------------------------------------------------
\section{CRM et gestion des interactions}
%--------------------------------------------------------------

\todo{%
  Redaction en PROSE uniquement (pas de tableau -- 3 alternatives sur des
  criteres simples ne justifient pas un tableau formel) : comparer EspoCRM,
  SuiteCRM, Odoo en 1-2 paragraphes (legerete, API REST, auto-hebergement,
  cout de licence). Conclure sur EspoCRM.
}

%--------------------------------------------------------------
\section{Reconnaissance émotionnelle}
%--------------------------------------------------------------

\todo{%
  Redaction en prose : justification metier (detecter detresse/agressivite
  pour escalade humaine, standard call center). Comparer approches audio
  (spectral) vs texte (NLP) -- texte retenu, justifier pourquoi (pas d'acces
  audio post-ASR, CPU contraint). Comparer les 3 modeles candidats
  (RoBERTa-large, XLM-R, DistilBERT multilingue) en 1 paragraphe avec
  tailles memoire et latence CPU -- un petit tableau 3 lignes est acceptable
  ici si le texte de justification reste developpe.
}

%--------------------------------------------------------------
\section{Portail de gouvernance — stack web}
%--------------------------------------------------------------

\todo{Backend : Django vs FastAPI vs Spring Boot.
Frontend : Vue vs React SPA vs Next.js.
Auth : Session vs JWT + 2FA.}

%--------------------------------------------------------------
\section{Tableau récapitulatif des choix technologiques}
%--------------------------------------------------------------

\begin{table}[H]
  \centering
  \caption{Récapitulatif des choix technologiques}
  \label{tab:recap_techno}
  \begin{tabularx}{\textwidth}{lXlX}
    \toprule
    \textbf{Couche} & \textbf{Évalué} & \textbf{Retenu} & \textbf{Critère décisif} \\
    \midrule
    Serveur SIP    & FreeSWITCH, Asterisk, Twilio & FreeSWITCH 1.10.11 & On-premise + mod\_audio\_stream \\
    VAD            & Silero V5, WebRTC VAD & Silero VAD V5 & Précision, Python natif \\
    STT            & Azure, Google, Whisper & faster-Whisper & Local, French fine-tuned \\
    NLU/Dialogue   & Rasa, Dialogflow, LUIS & Rasa 3.6.x & On-premise obligatoire \\
    Vector DB      & Pinecone, Weaviate, ChromaDB & ChromaDB & Local, léger \\
    Embeddings     & text-ada, bge-m3 & bge-m3 via Ollama & Multilingue, local \\
    LLM            & GPT-4, Llama, Cerebras & Cerebras gpt-oss-120b & Perf + conformité \\
    TTS            & ElevenLabs, Coqui, Piper & Piper + dict. phonémique & CPU, latence, local \\
    CRM            & Odoo, SuiteCRM, EspoCRM & EspoCRM & API REST, légèreté \\
    Émotion        & RoBERTa, XLM-R, DistilBERT & lxyuan/DistilBERT & CPU, multilingue, async \\
    Backend portail & Spring Boot, FastAPI, Django & Django + Celery & ORM, Celery, 2FA \\
    Frontend portail & Vue, React SPA, Next.js & Next.js & SSR, middleware auth \\
    \bottomrule
  \end{tabularx}
\end{table}

%--------------------------------------------------------------
\section{Conclusion}
%--------------------------------------------------------------

\todo{%
  Synthèse des choix justifiés par la contrainte de résidence des données.
  Transition : \og{}Les choix technologiques étant justifiés, le chapitre
  suivant présente la conception globale du système et l'articulation de
  ses 7 modules.\fg{}
}
